% Copyright (c) 2026 Sergey Zhunkov.
% SPDX-License-Identifier: CC-BY-4.0
%
% Set \ChartLanguage to en or ru before loading this source. English is the default.
\documentclass[a4paper,landscape]{article}
\usepackage[margin=0pt]{geometry}
\usepackage{fontspec,babel,etoolbox,tikz,luacode,hyperref,bookmark,lilyglyphs}
\usetikzlibrary{babel}

\providecommand{\ChartLanguage}{en}
\newif\ifChartRussian
\ifdefstring{\ChartLanguage}{ru}{%
  \ChartRussiantrue
  \babelprovide[import,main]{russian}
  \babelprovide[import]{english}
}{%
  \ifdefstring{\ChartLanguage}{en}{%
    \babelprovide[import,main]{english}
    \babelprovide[import]{russian}
  }{%
    \PackageError{transposition}{Unsupported ChartLanguage '\ChartLanguage'}%
      {Set ChartLanguage to en or ru.}%
  }%
}
\newcommand{\Localize}[2]{\ifChartRussian#1\else#2\fi}

\setmainfont{DejaVuSansCondensed.ttf}[BoldFont=DejaVuSansCondensed-Bold.ttf]
\pagestyle{empty}
\definecolor{ChartInk}{RGB}{23,43,50}
\definecolor{ChartMuted}{RGB}{70,91,99}
\definecolor{ChartRule}{RGB}{197,208,209}
\definecolor{ChartStripe}{RGB}{244,247,247}
\providecommand{\ChartAccentHex}{195755}
\definecolor{ChartAccent}{HTML}{\ChartAccentHex}
\colorlet{ChartHeader}{ChartAccent!13!white}
\colorlet{ChartHeading}{ChartAccent!33!ChartInk}

\newcommand{\ChartTitle}{\Localize{Транспонирование}{Transposition}}
\hypersetup{pdfauthor={Sergey Zhunkov},pdftitle={\ChartTitle},
  colorlinks=true,urlcolor=ChartHeading,linkcolor=ChartHeading}
\ifChartRussian
  \hypersetup{pdflang={ru-RU}}
\else
  \hypersetup{pdflang={en}}
\fi

\begin{luacode*}
local pitch_classes = {"C", "D♭", "D", "E♭", "E", "F", "F♯", "G", "A♭", "A", "B♭", "B"}
local alternate_names = {[2]="C♯", [4]="D♯", [7]="G♭", [9]="G♯", [11]="A♯"}

function PrintRootLookup()
  for column, note in ipairs(pitch_classes) do
    local alternate = alternate_names[column]
    tex.sprint(string.format([[\LookupRoot{%d}{%s}]], column,
      alternate and note .. "/" .. alternate or note))
  end
  for offset = 0, #pitch_classes - 1 do
    for column = 1, #pitch_classes do
      local target = pitch_classes[(column - 1 + offset) % #pitch_classes + 1]
      tex.sprint(string.format([[\LookupCell{%d}{%d}{%s}]], offset, column, target))
    end
  end
end
\end{luacode*}

% Coordinates and font sizes are in PDF points, measured from the bottom left.
\newcommand{\ChartText}[6][base west]{%
  \node[anchor=#1,inner sep=0pt,outer sep=0pt,text=#5,
    font=\fontsize{#4bp}{#4bp}\selectfont] at ({#2},{#3}) {#6};%
}
\newcommand{\ChartParagraph}[4]{%
  \node[anchor=north west,inner sep=0pt,outer sep=0pt,text=ChartInk,align=left,
    text width=#3bp,font=\fontsize{11bp}{14bp}\selectfont] at ({#1},{#2}) {#4};%
}
\newcommand{\ChartRuleAt}[1]{\draw[ChartRule,line width=.65bp] (34,#1) -- (808,#1);}
\newcommand{\ChartHeadingAt}[2]{\ChartText{34}{#1}{15}{ChartHeading}{\bfseries #2}}
\newcommand{\Card}[4]{%
  \draw[ChartRule,line width=.65bp,rounded corners=2bp] (#1,#2) rectangle (#3,#4);%
}
\newcommand{\ChartArrow}{→}
\newcommand{\ReferenceLink}[2]{%
  \href{https://github.com/zhunkoff/music-reference/blob/main/\Localize{README.ru.md}{README.md}\#ref-#1}{#2}%
}
\newcommand{\ChartPage}[3]{%
  \pdfbookmark[0]{#2}{transposition-\thepage}%
  \AddToHookNext{shipout/background}{%
    \put(0,-\paperheight){\begin{tikzpicture}[overlay,x=1bp,y=1bp]
      \ChartText{34}{542}{28}{ChartInk}{\bfseries\ChartTitle}%
      \ChartText[base east]{808}{542}{12}{ChartHeading}{\bfseries #1 / 2 · #2}%
      \ChartRuleAt{529}%
      #3%
      \ChartRuleAt{43}%
      \ChartText{34}{27}{9.5}{ChartMuted}{%
        \textbf{\Localize{См. также}{See also}}\quad
        \ReferenceLink{intervals}{\Localize{Интервалы}{Intervals}} ·
        \ReferenceLink{rhythm-and-meter}{\Localize{Ритм и размер}{Rhythm and meter}} ·
        \ReferenceLink{chord-construction}{\Localize{Построение аккордов}{Chord construction}} ·
        \ReferenceLink{chords-in-keys}{\Localize{Аккорды в тональностях}{Chords in keys}} ·
        \ReferenceLink{circle-of-fifths}{\Localize{Квинтовый круг}{Circle of fifths}}}%
      \ChartText[base east]{808}{27}{9}{ChartMuted}{%
        © 2026 \href{https://github.com/zhunkoff/music-reference}{Sergey Zhunkov} ·
        \href{https://creativecommons.org/licenses/by/4.0/}{CC BY 4.0}}%
    \end{tikzpicture}}%
  }%
  \null\newpage
}

\newcommand{\MethodCard}[3]{%
  \Card{#1}{371}{#1+250}{510}%
  \ChartText{#1+12}{489}{12.5}{ChartHeading}{\bfseries #2}%
  \ChartParagraph{#1+12}{475}{226}{#3}%
}
\newcommand{\LookupX}[1]{98+59.1667*(#1-.5)}
\newcommand{\LookupY}[1]{295-16*(#1+.5)}
\newcommand{\LookupRoot}[2]{\ChartText[base]{\LookupX{#1}}{301}{10.5}{ChartHeading}{\bfseries #2}}
\newcommand{\LookupCell}[3]{\ChartText[base]{\LookupX{#2}}{\LookupY{#1}-3.5}{11}{ChartInk}{#3}}
\newcommand{\RootLookup}{%
  \fill[ChartHeader] (34,295) rectangle (808,315);%
  \foreach \offset in {0,...,11}{%
    \ifodd\offset\relax
      \fill[ChartStripe] (34,{\LookupY{\offset}-8}) rectangle (808,{\LookupY{\offset}+8});%
    \fi
    \ChartText[base]{66}{\LookupY{\offset}-3.5}{10.5}{ChartHeading}{\bfseries +\offset}%
  }%
  \draw[ChartRule,line width=.65bp] (98,103) -- (98,315) (34,295) -- (808,295);%
  \Card{34}{103}{808}{315}%
  \ChartText[base]{66}{301}{10.5}{ChartHeading}{\bfseries\Localize{Сдвиг}{Shift}}%
  \directlua{PrintRootLookup()}%
}

\newcommand{\ScoreGlyph}[3]{\ChartText[base]{#1}{#2}{28}{black}{\lilyGlyph{#3}}}
\csdef{ScoreSignsharp}{accidentals.sharp}
\csdef{ScoreSignnatural}{accidentals.natural}
\newcommand{\ScoreQuarterStem}[2]{%
  \ifdim#2bp<14bp
    \draw[black,line width=.82bp] (#1+3.8,#2+.8) -- (#1+3.8,#2+24.5);%
  \else
    \draw[black,line width=.82bp] (#1-3.8,#2-.8) -- (#1-3.8,#2-24.5);%
  \fi
}
\newcommand{\MelodyX}[1]{\ifcase#1\or92\or145\or169\or231\or284\or308\fi}
\newcommand{\MelodyScore}[4]{%
  \ChartText{#1}{468}{12}{ChartHeading}{\bfseries #2}%
  \begin{scope}[shift={(#1,420)}]
    \foreach \liney in {0,7,14,21,28}{\draw[black,line width=.46bp] (0,\liney) -- (356,\liney);}%
    \draw[black,line width=.65bp] (0,0) -- (0,28) (205,0) -- (205,28) (351,0) -- (351,28);%
    \draw[black,line width=2bp] (355,0) -- (355,28);%
    \ScoreGlyph{12}{7}{clefs.G}%
    #3%
    \ScoreGlyph{70}{14}{two}%
    \ScoreGlyph{70}{0}{four}%
    % Each bar uses one quarter and a beamed pair of eighth notes.
    \foreach \step/\pitch/\sign [count=\event] in {#4}{%
      \pgfmathsetmacro{\NoteX}{\MelodyX{\event}}%
      \pgfmathsetmacro{\NoteY}{3.5*\step}%
      \ifnum\step=-2\relax
        \draw[black,line width=.6bp] (\NoteX-6,\NoteY) -- (\NoteX+6,\NoteY);%
      \fi
      \ifdefstring{\sign}{none}{}{\ScoreGlyph{\NoteX-10}{\NoteY}{\csuse{ScoreSign\sign}}}%
      \ScoreGlyph{\NoteX}{\NoteY}{noteheads.s2}%
      \ifnum\event=1\relax
        \ScoreQuarterStem{\NoteX}{\NoteY}%
      \else\ifnum\event=4\relax
        \ScoreQuarterStem{\NoteX}{\NoteY}%
      \else
        \draw[black,line width=.82bp] (\NoteX+3.8,\NoteY+.8) -- (\NoteX+3.8,35);%
      \fi\fi
      \ChartText[base]{\NoteX}{-22}{10}{ChartMuted}{\pitch}%
    }%
    \draw[black,line width=2.7bp] (148.8,34) -- (172.8,34) (287.8,34) -- (311.8,34);%
  \end{scope}%
}
\newcommand{\EFlatSignature}{%
  \ScoreGlyph{30}{14}{accidentals.flat}%
  \ScoreGlyph{41}{24.5}{accidentals.flat}%
  \ScoreGlyph{52}{10.5}{accidentals.flat}%
}
\newcommand{\ProgressionRow}[6]{%
  \ChartText{46}{#1}{11}{ChartHeading}{\bfseries #2}%
  \ChartText[base]{395}{#1}{12}{ChartInk}{#3}%
  \ChartText[base]{508}{#1}{12}{ChartInk}{#4}%
  \ChartText[base]{621}{#1}{12}{ChartInk}{#5}%
  \ChartText[base]{734}{#1}{12}{ChartInk}{#6}%
}

\begin{document}
\ChartPage{1}{\Localize{Метод и таблица}{Method and lookup}}{%
  \MethodCard{34}{\Localize{1. Задайте перенос}{1. Set the transposition}}{\Localize
    {Выберите исходную и целевую тональности одного лада. Задайте
     направление и регистр. Из до мажора в ми-бемоль мажор перенос
     C4~\ChartArrow~E♭4 даёт малую терцию вверх (+3 полутона).
     Примените тот же сдвиг ко всем звукам.}
    {Choose the source and target keys in the same mode. Set the direction
     and register. From C major to E♭ major, C4~\ChartArrow~E♭4 moves up
     a minor third (+3 semitones). Apply the same shift to every pitch.}}%
  \MethodCard{296}{\Localize{2. Отсчитайте буквы и полутоны}{2. Count letters and semitones}}{\Localize
    {При переносе по интервалу отсчитайте буквы в нужном направлении,
     считая исходную за первую. Подберите знак по числу полутонов.
     C~\ChartArrow~D♭ и C~\ChartArrow~C♯ дают +1. Первая пара образует
     малую секунду, вторая — увеличенную приму.}
    {For a named interval, count letters in the chosen direction,
     starting with the original letter as one. Choose the accidental
     for the semitone distance. C~\ChartArrow~D♭ and C~\ChartArrow~C♯
     both give +1. The first is a minor second, the second an augmented unison.}}%
  \MethodCard{558}{\Localize{3. Проверьте запись и регистр}{3. Check spelling and register}}{\Localize
    {Прибавьте сдвиг в полутонах со знаком. Выберите написание по целевой
     тональности и роли ноты. Номер октавы меняется на C.
     На полутон вверх B4~\ChartArrow~C5, вниз C4~\ChartArrow~B3.
     Сдвиг ±12 сохраняет класс высоты и меняет октаву.}
    {Add the signed semitone count. Choose the spelling for the target
     key and the note’s role. Octave numbers change at C. Up one semitone,
     B4~\ChartArrow~C5. Down one, C4~\ChartArrow~B3. A shift of ±12
     keeps the pitch class and changes the octave.}}%

  \ChartHeadingAt{352}{\Localize{Хроматическая таблица основных тонов}{Chromatic root lookup}}%
  \ChartParagraph{34}{343}{774}{\Localize
    {Выберите исходный основной тон вверху и строку сдвига в полутонах. Строки показывают подъём в пределах октавы.
     Для движения вниз прибавляйте 12 к отрицательному сдвигу, пока он не станет неотрицательным.}
    {Find the source root along the top and the semitone shift on the left. Rows go up within an octave.
     For downward shifts, add 12 to the negative shift until it is nonnegative.}}%
  \RootLookup
  \ChartParagraph{34}{92}{774}{\Localize
    {Например, −3 даёт строку +9. Из E♭ получаем C. Для сдвигов от +12 и больше вычитайте 12 до попадания в диапазон 0–11.
     В ячейках по одному имени на класс высоты. Написание выбирайте по контексту, а октаву считайте отдельно.
     Названия международные, от C до B (B = си).}
    {For example, −3 uses row +9. From E♭, that gives C. For shifts of +12 or more, subtract 12 until the result is in 0–11.
     Cells show one name per pitch class. Choose the spelling in context and work out the octave separately.
     Note names use the international C–B system.}}%
}

\ChartPage{2}{\Localize{Мелодия и аккорды}{Melody and chords}}{%
  \ChartHeadingAt{505}{\Localize{Мелодия · на малую терцию вверх (+3)}{Melody · up a minor third (+3)}}%
  \Card{34}{386}{808}{485}%
  \MelodyScore{46}{\Localize{Исходная · до мажор}{Source · C major}}{}%
    {-2/C4/none,0/E4/none,1/F♯4/sharp,2/G4/none,0/E4/none,-2/C4/none}%
  \MelodyScore{440}{\Localize{Целевая · ми-бемоль мажор}{Target · E♭ major}}{\EFlatSignature}%
    {0/E♭4/none,2/G4/none,3/A♮4/natural,4/B♭4/none,2/G4/none,0/E♭4/none}%
  \ChartParagraph{34}{376}{774}{\Localize
    {Запишите ключевые знаки целевой тональности. F♯ переходит в A♮, повышенную IV ступень.
     Бекар отменяет ключевой бемоль у A, а E♭ и B♭ получают бемоли по ключу.
     Интервалы мелодии и ритм сохраняются.}
    {Write the target key signature. F♯ becomes A♮, the raised fourth degree.
     The natural cancels the signature’s A♭. E♭ and B♭ take their flats from the signature.
     The melody keeps its intervals and rhythm.}}%

  \ChartHeadingAt{333}{\Localize{Аккордовые последовательности}{Chord progressions}}%
  \ChartParagraph{34}{322}{774}{\Localize
    {Римскими цифрами обозначьте ступени основных тонов. В целевой тональности найдите те же ступени
     и перенесите качество аккордов и надстройки. Прописные цифры обозначают мажорное трезвучие, строчные — минорное.}
    {Use Roman numerals for the roots’ scale degrees. Find the same degrees in the target key and carry over
     each chord’s quality and extensions. The numerals are uppercase for major triads and lowercase for minor triads.}}%
  \fill[ChartHeader] (34,274) rectangle (808,290);%
  \fill[ChartHeader] (34,225) rectangle (808,241);%
  \Card{34}{192}{808}{290}%
  \ProgressionRow{278}{\Localize{Большая секунда вверх (+2)}{Up a major second (+2)}}{I}{ii}{V7}{I}%
  \ProgressionRow{262}{\Localize{Исходная · до мажор}{Source · C major}}{C}{Dm}{G7}{C}%
  \ProgressionRow{246}{\Localize{Целевая · ре мажор}{Target · D major}}{D}{Em}{A7}{D}%
  \ProgressionRow{229}{\Localize{Большая секунда вниз (−2)}{Down a major second (−2)}}{i}{iv}{V7}{i}%
  \ProgressionRow{213}{\Localize{Исходная · ля минор}{Source · A minor}}{Am}{Dm}{E7}{Am}%
  \ProgressionRow{197}{\Localize{Целевая · соль минор}{Target · G minor}}{Gm}{Cm}{D7}{Gm}%
  \ChartText{34}{178}{10.5}{ChartInk}{\Localize
    {В минорной последовательности повышенная VII ступень в V7 переносится из G♯ в F♯.}
    {In the minor progression, the raised seventh degree in V7 moves from G♯ to F♯.}}%

  \Card{34}{68}{528}{165}%
  \ChartText{46}{148}{12.5}{ChartHeading}{\bfseries\Localize{Бас, надстройки и альтерации}{Bass, extensions, and alterations}}%
  \ChartParagraph{46}{136}{470}{\Localize
    {Из до мажора в ре мажор (+2) C/E переходит в D/F♯. Основной тон
     C~\ChartArrow~D и бас E~\ChartArrow~F♯ сдвигаются вместе, поэтому терция
     остаётся в басу. G7(♭9) переходит в A7(♭9). Септима остаётся малой,
     а пониженная нона A♭ становится B♭.}
    {From C major to D major (+2), C/E becomes D/F♯. The root C~\ChartArrow~D and bass
     E~\ChartArrow~F♯ move together, keeping the third in the bass.
     G7(♭9) becomes A7(♭9). The seventh stays minor and the flatted ninth
     A♭ becomes B♭.}}%
  \Card{540}{68}{808}{165}%
  \ChartText{552}{148}{12.5}{ChartHeading}{\bfseries\Localize{Транспонирование и каподастр}{Transposition and capo}}%
  \ChartParagraph{552}{136}{244}{\Localize
    {Полный каподастр повышает звучание тех же гитарных аппликатур
     на полутон за лад. Общий метод применим к нотам и аккордам
     для любого инструмента.}
    {A full guitar capo raises the sound of the same shapes by a semitone
     per fret. The general method applies to written notes and chords
     for any instrument.}}%
  \ChartText{34}{55}{9.5}{ChartMuted}{%
    \textbf{\Localize{Источники}{Sources}}\quad
    Steps to Music Theory ·
    \href{https://odp.library.tamu.edu/stepstomusictheory/chapter/more-intervals/}
      {\Localize{перенос по интервалу}{intervallic transposition}}\qquad
    Open Music Theory ·
    \href{https://pressbooks.nebraska.edu/openmusictheory/chapter/chord-symbols/}
      {\Localize{символы аккордов}{chord symbols}}\qquad
    University of Puget Sound ·
    \href{https://musictheory.pugetsound.edu/mt21c/RomanNumeralChordSymbols.html}
      {\Localize{римские цифры}{Roman numerals}}}%
}
\end{document}
